\subsection{terminal}\label{terminal}
The \verb|terminal| module is a \verb|kybd_t| implementation (Section
\ref{keyboard}) that accepts key input from a serial console. It exports the
device \verb|terminal_dev| with the operating mode \verb|TERMINAL|. Internally
it keeps its own \verb|state_t| (\verb|my_term|) with eight keys labelled
\verb|'1'|\dots\verb|'8'|, and for UART access it relies on the
\verb|sio_t| configuration of the \verb|serial| module (Section \ref{serial}).

\subsubsection*{Operation}
{\sloppy\emergencystretch=3em
On scanning (\verb|terminal_scan|) the module prompts once for input
(``Please enter key'') and then reads character by character via
\verb|HAL_UART_Receive| up to the line ending. A recognised digit is mapped to
a key index via \verb|state_ch2idx()|; if it lies within the valid range, the
change is propagated into \verb|my_term| and the \emph{dirty} flag is set. The
remaining vtable methods (\verb|terminal_state|, \verb|terminal_diff|,
\verb|terminal_add|, \verb|terminal_isdirty|, \dots) access the functions of
the \verb|state| module directly.
\par}

\subsubsection*{API}
{\sloppy\emergencystretch=3em
\begin{description}
\item[\texttt{terminal\_waitForNumber(key)}] Blocking: reads a number from the
  serial console; \verb|key| then points to the read string. The input
  \verb|R|/\verb|r| signals a reset. Returns: \verb|int8_t| -- the value, or
  $-1$ on reset/invalid input.
\end{description}
\par}

\noindent In addition there is a \verb|UART_IdleCallback()| for DMA-based
reception (idle-line detection), which takes received characters into the
label \verb|clabel| of the \verb|state_t|.
